\documentclass[a4paper,10pt]{ltjsarticle}
\usepackage[margin=22mm]{geometry}
\usepackage{amsmath,booktabs}
\title{回転補償ホログラムに関する疑問点}
\author{}
\date{2026年11月27日}
\begin{document}
\maketitle
\section{現時点で確認すべき質問}
\subsection*{① 初期傾斜角を0度として扱ってよいのか}
傾斜ポリゴンから生成した複素ホログラムは、グローバル座標系の $z_0=0$ 面に定義される。この記録面の姿勢を根拠に変換式で $R_1=I$ としてよいか。物体の初期傾斜角は、入力波面と変換式のどこに反映されるか。

\subsection*{② 光波のキャリアの基準方向は適切か}
従来の導出では、$P_1\to P_0$ に $z_1$ 軸方向、$P_0\to P_2$ に $z_2$ 軸方向のキャリアを用いた。ホログラム面を固定する今回の目的には、入出力を共通の $z_0$ 軸を基準に定義する方が適切か。

\subsection*{③ キャリアシフトの扱いは適切か}
直接変換式には、局所キャリアに由来する横周波数シフトの差
\[
  (b_3-a_3)/\lambda,\qquad (b_6-a_6)/\lambda
\]
が含まれる。$R_1=I$ では $a_3=a_6=0$ となる。この入力側のキャリア規約は、傾斜ポリゴンから生成された複素ホログラムの位相の定義とも整合するか。

\subsection*{④ 傾斜平面間変換と物体回転は等価か}
導出した理論は、異なる傾斜平面間の角スペクトル変換である。研究目的は、ホログラム面を固定したまま物体を回転させた場合のホログラムを得ることである。照明、キャリアの方向、回転中心、伝搬条件を含め、どの条件で両者が等価になるか。

\subsection*{⑤ 相対回転行列だけで変換式を表せるか}
現在の式には $R_1$ と $R_2$ の成分 $a_i,b_i$ がそれぞれ残る。
$R_\Delta=R_2R_1^{\mathsf T}$ を用いて、キャリアシフトと縦周波数の平方根の項を含めた写像を表せるか。その際、共通基準軸に対する各面の姿勢など、相対回転以外に必要な情報は何か。

\subsection*{⑥ 二つの変換理論を統合した目的は何か}
常に $R_1=I$ とする条件では、$P_1\to P_0\to P_2$ の座標写像は $P_0\to P_2$ に帰着する。$P_1\to P_0$ を含む一般式は、どの実験条件や将来の用途で必要になるか。

\subsection*{⑧ FFTスペクトルの傾きが一致しないのはなぜか}
45度のOriginalから30度を公称出力角とする変換を行うと、30度のOriginalと変換結果で、FFT振幅スペクトルの斜め成分の傾きが異なる。補間に由来する差と、参照周波数座標の写像に由来する差をどのように切り分けるか。

\subsection*{⑨ 三角形の横幅が一致しないのはなぜか}
初期物体角を $\theta_0$、相対回転角を $\delta$ とすると、幾何学的な投影幅の変換率と、現在の周波数座標写像を原点近傍で一次近似した変換率は、それぞれ
\[
 M_{\mathrm{original}}=\frac{\cos(\theta_0+\delta)}{\cos\theta_0},
 \qquad
 M_{\mathrm{transform}}\simeq\frac{1}{\cos\delta}
\]
である。45度から30度では、前者は $1.2247$、後者は $1.0353$ となる。
この差は、三角形の形状やFFTの斜め成分の傾きの違いを説明できるか。
また、変換に初期物体角が陽に入っていないこととどのように関係するか。
現段階では一次近似に基づく仮説であり、厳密な証明と実測幅との照合は今後行う。

\section{今後の研究課題}
\subsection*{⑫ 初期物体角を明示せずに回転変換できるか}
計算済みの複素ホログラムと相対角 $\delta$ だけから、任意の初期物体角 $\theta_0$ に対して
\[
 H_{\theta_0+\delta}=T_\delta[H_{\theta_0}]
\]
を実現できるか。入力波面に含まれる情報で足りる条件と、物体姿勢などの追加情報が必要な条件を明確にする。

\subsection*{⑬ 複数ポリゴンを一括で回転変換できるか}
複数ポリゴンの合成ホログラムを $H=\sum_{i=1}^{N}H_i$ とする。
各成分に共通の線形演算子 $T_\delta$ を適用できるなら、
\[
 T_\delta[H]
 =T_\delta\!\left[\sum_{i=1}^{N}H_i\right]
 =\sum_{i=1}^{N}T_\delta[H_i]
\]
となる。各ポリゴンの頂点、初期傾斜角、奥行きを再参照せず、この演算子で回転後の各成分を再現できる条件は何か。遮蔽の変化は別の発展課題として扱う。

\section{確認の順序}
\begin{center}
\begin{tabular}{@{}lll@{}}
\toprule
分類 & 番号 & 確認項目\\
\midrule
理論の前提 & ①--⑥ & 初期姿勢、キャリア、変換の等価性、一般式の役割\\
実験との照合 & ⑧・⑨ & FFTの斜め成分と三角形の横幅\\
今後の研究 & ⑫・⑬ & 初期物体角を明示しない変換と複数ポリゴンへの適用\\
\bottomrule
\end{tabular}
\end{center}
まず理論の前提を確認し、FFTと横幅の結果との整合性を検証する。そのうえで、初期物体角を明示しない変換と複数ポリゴンへの適用を検討する。
\end{document}
